\documentclass{article}
\usepackage[utf8]{inputenc} 
\usepackage[spanish]{babel} 
\usepackage{tikz}
\usepackage{graphicx} 
\usepackage{listings}
\usepackage{xcolor}
\definecolor{codegreen}{rgb}{0,0.6,0}
\definecolor{codegray}{rgb}{0.5,0.5,0.5}
\definecolor{codepurple}{rgb}{0.58,0,0.82}
\definecolor{backcolour}{rgb}{0.95,0.95,0.92}
\lstset{
    backgroundcolor=\color{backcolour},   
    commentstyle=\color{codegreen},
    keywordstyle=\color{magenta},
    numberstyle=\tiny\color{codegray},
    stringstyle=\color{codepurple},
    basicstyle=\ttfamily\footnotesize,
    breaklines=true,                 
    captionpos=b,                    
    keepspaces=true,                 
    numbers=left,                    
    showspaces=false,                
    language=C++
}
\begin{document}
\begin{center}
    \huge \textbf{Estructura de Datos: DSU}
    \vspace{0.5cm}
\end{center}
La estructura \emph{Disjoint Set Union} (DSU), también llamada \emph{Union-Find}, permite gestionar una partición de un conjunto de elementos en varios conjuntos disjuntos. Cada conjunto está identificado por un \textbf{representante}, que suele ser el nodo raíz del árbol que los agrupa.

Sea $V$ un conjunto de tamaño $n$, inicialmente cada elemento $v_i$ es el representante de su propio conjunto de un solo elemento:
\begin{center}
    $V=\{v_1, v_2, v_3, \dots, v_n\}$    
\end{center}



Esta estructura destaca por dos operaciones altamente eficientes:
\begin{itemize}
    \item \textbf{Find}: Determina a qué conjunto pertenece un elemento devolviendo su representante.
    \item \textbf{Union}: Fusiona dos conjuntos distintos en uno solo.
\end{itemize}

\section{Construcción e Inicialización}
Para implementar DSU, utilizamos un vector $P$ para almacenar el padre de cada nodo y un vector $SZ$ para registrar el tamaño de cada conjunto. Inicialmente, cada nodo es su propio padre ($P[i] = i$) y el tamaño de cada conjunto es 1.

\begin{lstlisting}[language=C++, caption=Constructor de la clase DSU]
vector<int> p, sz;
void build(int n) {
    p.resize(n + 1);
    sz.assign(n + 1, 1);
    for(int i = 1; i <= n; i++) p[i] = i;
}
\end{lstlisting}

\section{Find}
Para hallar la raíz de un nodo $i$, ascendemos recursivamente por sus padres hasta encontrar un nodo tal que $P[i] == i$. 
\textbf{Optimización (Path Compression):} Durante la recursión, actualizamos $P[i]$ directamente con el valor de la raíz hallada. Esto "aplana" el árbol, acelerando futuras consultas de forma drástica.

\begin{lstlisting}[language=c++, caption=Funcion Find con Compresion de Caminos]
int find(int i){
    if(p[i]==i)return i;
    return p[i]=find(p[i]); 
}   
\end{lstlisting}



\section{Unite}
Al unir dos conjuntos, primero verificamos que no tengan ya el mismo representante. 
\textbf{Optimización (Union by Size):} Conectamos siempre el árbol más pequeño bajo la raíz del árbol más grande. Esto mantiene la altura del árbol mínima ($O(\log n)$), evitando que la estructura se degrade.

\begin{lstlisting}[language=c++, caption=Funcion Union por Tamaño]
void unite(int i, int j) {
    int ri=find(i);
    int rj=find(j);
    if(ri!=rj) {
        if(sz[ri]<sz[rj])swap(ri,rj);
        p[rj]=ri;
        sz[ri]+=sz[rj];
    }
}
\end{lstlisting}
\end{document}